\section{Conclusions}\label{sec:conclusions}

Across many education systems, peers have been shown to influence a student's own academic achievement (\cite{sacerdote2011peer}). Understanding the mechanisms behind this influence could shed light on how school environments shape early differences in achievement, which persist over time with major lifelong consequences (\cite*{cunha2006interpreting,heckman2014economics}). But empirical challenges have hindered progress towards this goal (\cite*{blume2011identification}). This article exploits a new empirical context with rich data on how students and schools responded to earthquake-induced study disruptions to examine how disruptions spill over to peers' learning and develop a new theory of peer influence in schools.\par 


Exploiting the context of one of the most violent earthquakes ever recorded and detailed data on the disruptions to each student's home environment, the study finds that disruptions can lower the reported ease of exerting study effort, with negative consequences for achievement that persist for at least 22 months. Notably, such disruptions can spill over to classmates, affecting their achievement.\footnote{There is evidence that environmental risks can spill over to classmates also in the context of lead exposure (\cite{gazze2021long}). } Schools mediated some, but likely not all, of these spillover effects, suggesting a possible mediating role for peer interactions. Following \cite*{blume2015linear}, I micro-found the observed spillovers through a model of student interactions, that also allows for school mitigation, and I derive comparative statics that rationalize the empirical findings. I show that the empirical evidence is consistent with a mode of interaction that has not received much attention in the peer effects literature before: competition for classroom rank. A desire to compete implies that moments beyond the mean of peer characteristics matters, which is an empirical fact across several settings.\footnote{In contrast, a desire-to-conform assumption underlies empirical identification strategies that contrast within- and across-group variances in outcomes to identify excess variance across groups that cannot be explained by individual and group heterogeneity and/or selection (\cite{graham2008identifying} and \cite{glaeser1996crime}). Whenever a desire to compete is the true interaction mode, such methods may fail to detect peer effects when they are present, a false negative result.   }\par 









The results offer new insights for policy. Peer assignment policies, such as tracking students by ability, are among the most commonly studied in the schooling context (e.g. \cite{duflo2011peer,garlick2018academic}). My results suggest that their impacts could vary depending on whether performance rank is intrinsically or extrinsically rewarded. Ability tracking may improve the achievement of all students, even those in the lower-tracks, in settings in which students care about their performance rank, by increasing the number of nearby competitors. There are several reasons why students may value rank: competitive preferences emerge early on, and a growing body of evidence shows that class rank offers immediate and long-term benefits. In many education settings, intrinsic rank concerns are reinforced by explicit rank-based incentives. Teachers often grade on a curve (\cite{calsamiglia2019grading}), and higher education admissions frequently rely on within-school rankings (\cite{horn2003percent,grau2018impact,carlana2024far,tincani2025college}). The mechanism identified in this paper, therefore, may operate broadly.\par 

Much is still unknown about the interaction between rank-based rewards and classroom allocation rules. Measuring intrinsic rank concerns and extrinsic rank rewards in schools could become a way to inform the targeting of grouping policies. Future research could also compare the achievement gains from optimally designing rank rewards and group allocations to the potential labor market losses from lower prosociality due to enhanced competition (\cite{kosse2020prosociality,chen2022competition,kosse2023ps}). Answering these open questions could significantly advance our understanding of social interactions in school, and expand our toolkit of cost-effective policy interventions.\par 
 



